\documentclass[../../main.tex]{subfiles}
\begin{document}
\section{Family 5: transport maps and the entropic barrier}
\label{sec:family-transport}

\paragraph{Object followed.}
A map pushing a reference measure --- usually a Gaussian, or Wiener measure --- onto $\mu$,
together with the derivative of that map.

\paragraph{What it buys.}
A reformulation in which KLS becomes a regularity statement about a single transport map rather
than a spectral statement about a diffusion. The reformulation is exact and appealing; the
difficulty is that the three natural instances each fail in an instructive and different way.

\subsection{Caffarelli}
\label{subsec:transport-caffarelli}

If the target is uniformly more log-concave than a Gaussian, Caffarelli's contraction theorem makes
the Brenier map from that Gaussian to $\mu$ $1$-Lipschitz. The Gaussian Poincar\'e inequality then
transfers directly, with no loss.

This settles the strongly log-concave case, the same case already settled by Bakry--\'Emery for
structured families. It cannot cover arbitrary log-concave measures, whose curvature may
vanish identically --- the uniform measure on a convex body being the extreme instance.

\subsection{Brownian (F\"ollmer) transport}
\label{subsec:transport-brownian}

The Brownian transport map sends Wiener measure to a target $\mu$ \cite{MikulincerShenfeld2021BrownianTransport}.
A dimension-free estimate on its Malliavin derivative,
\begin{equation}
  \E\norm{\calD X_1}_\op^2\le C,
  \label{eq:brownian-derivative}
\end{equation}
would give $\Var_\mu f\le C\Lip(f)^2$ for every Lipschitz $f$; Milman's equivalence between
Lipschitz concentration and a spectral gap for log-concave measures
\cite{Milman2009Isoperimetric} would then upgrade that to KLS.

\begin{warning}[No pathwise dimension-free Lipschitz map exists]
\label{rem:no-lipschitz-transport}
The corresponding \emph{pathwise} statement is false, and cheaply so. A dimension-free Lipschitz
map from a Gaussian to every log-concave $\mu$ would transfer the Gaussian logarithmic Sobolev
inequality and Gaussian tail decay to $\mu$. Already the one-dimensional exponential law violates
both. So \eqref{eq:brownian-derivative} must be read as an \emph{averaged operator} estimate; the
$\E$ and the order of the norm are not cosmetic.
\end{warning}

This is the reason the family has not yet produced an independent bound: current estimates on
\eqref{eq:brownian-derivative} are polylogarithmic and largely inherit existing KLS-scale input
rather than supplying it. The averaged operator derivative remains the plausible target.

\subsection{The entropic barrier}
\label{subsec:transport-entropic}

For a convex body $K$ set
\[
  \Lambda(\theta)=\log\int_Ke^{\inner\theta x}\dd x .
\]
Then
\begin{equation}
  \Hess\Lambda(\theta)=\Cov_{p_\theta}(X),
  \label{eq:entropic-hessian}
\end{equation}
where $p_\theta\propto e^{\inner\theta x}\one_K$, and $D^3\Lambda(\theta)$ is the third-moment
tensor of the same exponential tilt \cite{BubeckEldan2014EntropicBarrier}. The entropic barrier
therefore packages exactly the covariance and third-moment geometry that stochastic localization
manipulates --- \eqref{eq:entropic-hessian} and \eqref{eq:sl-covariance-sde} are two descriptions
of one object, one indexed by a tilt parameter and one by a time.

The gap is a matter of which contractions are controlled. Ordinary self-concordance controls
\emph{scalar} contractions of $D^3\Lambda$, of the form $D^3\Lambda[u,u,u]$. Localization needs
\emph{matrix} or Hilbert--Schmidt contractions --- precisely the parameter $\kappa_n$ of
\eqref{eq:kappa-def}. Letwin's theorem now supplies the latter
(Proposition~\ref{prop:letwin-kappa}), which closes this particular mismatch; what it does not
supply is the all-functions spectral step, and that remains as it was.

\paragraph{The precise missing estimate.}
The dimension-free averaged bound \eqref{eq:brownian-derivative} on the expected operator norm of
the transport derivative.

\paragraph{Why it stalls.}
Warning~\ref{rem:no-lipschitz-transport} rules out the pathwise route, so any argument must work
in expectation --- and an expected operator norm is not accessible by the pathwise convexity
techniques that make Caffarelli's theorem work. Available derivative bounds either lose the
alignment between the derivative and the test function, or reuse KLS-scale input and so cannot
improve it.

\begin{remark}[Structured cases where KLS is known]
\label{rem:known-cases}
It is worth recording what is not open, since it delimits what a counterexample could look like.
KLS holds for products by tensorization; for uniformly log-concave measures via
Brascamp--Lieb/Bakry--\'Emery (Section~\ref{subsec:transport-caffarelli}); for $\ell_p$-balls; and
for broad classes of generalized Orlicz balls \cite{KolesnikovMilman2016OrliczKLS}. It is
\emph{not} known for all unconditional convex bodies. These cases illustrate how additional
structure can make the corresponding constant explicit; KLS asks for a universal bound when no
such structure is available.
\end{remark}

\paragraph{How this family feeds the live routes.}
It feeds none of them, and the gap it leaves is recorded as such. Target~4 of
Section~\ref{sec:kls-synthesis} --- extending parallel coupling beyond linear tilts --- has no
corresponding node in this repository's route registry, and the synthesis says so rather than
inventing one. Transport enters Parts~III and~IV only through the Brascamp--Lieb cap that
stochastic localization uses.
\end{document}
